\begin{table}[h]
\caption{Generality of magnitude trim across data-free merging methods. For each method, we apply identical TIES-style magnitude trim with keep fraction $k$ to each $\Delta_t$ before running the method (3-seed mean $\pm$ std, all 7 tasks). Reported is the best of $k\in\{0.3, 0.5\}$. Only ACTMat's covariance solve benefits substantially on both architectures; for the spectral method Iso-C the same trim \emph{hurts}. The TACT improvement is therefore specific to interference-minimizing covariance solves, not a generic side-effect of trimming.}
\label{tab:trim_other}
\centering
\small
\setlength{\tabcolsep}{4pt}
\begin{tabular}{lcc|cc}
\toprule
 & \multicolumn{2}{c|}{\textbf{ViT-B/32}} & \multicolumn{2}{c}{\textbf{ViT-B/16}} \\
Method & Baseline & Best ($+\Delta$) & Baseline & Best ($+\Delta$) \\
\midrule
Task Arithmetic & 72.99 & 74.10 ($+1.11$) & 78.82 & 80.11 ($+1.29$) \\
Iso-C & 77.40 & 77.05 ($-0.35$) & 83.97 & 83.69 ($-0.28$) \\
TSV & 73.00 & 73.93 ($+0.92$) & 78.89 & 79.99 ($+1.10$) \\
\textbf{ACTMat} & \textbf{83.86} & \textbf{86.13 (+2.28)} & \textbf{87.63} & \textbf{88.79 (+1.15)} \\
\bottomrule
\end{tabular}
\end{table}
